\section*{Nicht verf\"ugbare Angaben und Grenzen des Berichts}
\addcontentsline{toc}{section}{Nicht verf\"ugbare Angaben}
\markright{Nicht verf\"ugbare Angaben}

\begin{description}\small
\item[Marktanteile.] GEA nennt keine. Die Behauptung der F\"uhrungsposition
  (Abschnitt~\ref{sec:moat}) bleibt deshalb unbelegt; sie geht in keine Rechnung ein.
\item[EBITDA vor Restrukturierung und Marge f\"ur 2019.] Der Gesch\"aftsbericht 2020
  weist die Kennzahl f\"ur das Vorjahr nicht in vergleichbarer Abgrenzung aus; die Reihe
  beginnt bei diesen beiden Gr\"o{\ss}en deshalb 2020.
\item[Organisches Umsatzwachstum vor 2022.] Erst ab dem Gesch\"aftsbericht 2023 auf der
  Kennzahlenseite gedruckt. Die Prognosetreue f\"ur 2021 bleibt beim Umsatz dadurch
  unbewertet.
\item[Einzelne Kursziele der Analysten.] Nur Mittelwert, H\"ochst- und Tiefstwert liegen
  vor; die Verteilung dazwischen ist unbekannt.
\item[Leerverkaufsquote.] Nicht beschafft.
\item[Vollst\"andigkeit der Directors' Dealings.] Die Meldungen wurden einzeln abgerufen,
  nicht als amtlicher Registerexport. Weitere K\"aufe sind m\"oglich; weitere Verk\"aufe
  w\"urden das Bild \"andern.
\item[Widerspruch im Gesch\"aftsbericht 2025.] Die EBITDA-Marge 2025 ist mit
  \Pz{\MargeVergleich} und mit \Pz{\EbitdaMarge} zweimal unterschiedlich gedruckt. Der
  Bericht rechnet mit dem aus zwei anderen gedruckten Werten folgenden Wert
  (\MargeBerechnet) und h\"alt den Widerspruch fest, statt ihn aufzul\"osen.
\item[Grenze des Verfahrens.] Die Seitenpr\"ufung stellt sicher, dass jede Zahl auf der
  genannten Seite steht. Sie kann nicht ausschlie{\ss}en, dass eine dort stehende Zahl
  falsch zugeordnet wurde. Dagegen stehen die Kettenpr\"ufung \"uber zwei Berichte und
  der Abgleich von Kennzahlenseite und Kapitalflussrechnung, nicht mehr.
\end{description}
